\subsection{Cuadrante de tareas domésticas y asistente inteligente}
\label{subsec:tareas_asistente}

El último eslabón funcional de la plataforma aborda la organización del trabajo doméstico compartido y la integración de un asistente conversacional inteligente para dar soporte autónomo a la comunidad.

\subsubsection{Cuadrante rotativo de tareas y gamificación}

Para erradicar los desacuerdos derivados del reparto desigual de tareas de limpieza y mantenimiento, Colivi implementa en \texttt{ChoresPage.tsx} un planificador automatizado respaldado por un sistema de recompensas:

\begin{itemize}
    \item \textbf{Series recurrentes y rotación Round-Robin}: El administrador del hogar define las tareas periódicas (ejemplo: «Limpieza de cocina los domingos» o «Tirar la basura reciclable»). El motor del backend distribuye los turnos de forma secuencial y equitativa entre los miembros activos del piso mediante un algoritmo rotativo \emph{Round-Robin}, garantizando que la carga de trabajo se distribuya de manera perfectamente homogénea a lo largo del curso académico.
    \item \textbf{Mecánica gamificada de «Rescate de tareas»}: Durante periodos de alta exigencia académica (como épocas de exámenes finales), un estudiante puede encontrarse imposibilitado para ejecutar su turno asignado. En lugar de generar conflictos, el sistema permite activar un botón de «Solicitar rescate». Cualquier otro compañero de piso puede asumir la tarea voluntariamente a cambio de recibir puntos de karma de la comunidad, incentivando la cooperación solidaria.
    \item \textbf{Verificación y puntuación acumulativa}: Al completar una labor, el usuario la marca como realizada, actualizando su puntuación en la tabla clasificatoria del hogar y registrando la acción en el muro de trazabilidad común.
\end{itemize}

\subsubsection{Canal de mensajería contextual y copiloto inteligente con protocolo MCP}

La comunicación y el soporte se articulan en el frontend mediante la bandeja de mensajes (\texttt{MessagesPage.tsx}) y el widget del asistente inteligente:

\begin{enumerate}
    \item \textbf{Chat directo contextual}: Permite la conversación privada entre inquilinos y propietarios para pactar visitas, así como un canal grupal interno exclusivo para los miembros del hogar.
    \item \textbf{Asistente conversacional contextual}: Integrado en la interfaz mediante \textbf{Spring AI}, el usuario puede interactuar en lenguaje natural con un copiloto inteligente. Dicho asistente no se limita a responder preguntas estáticas, sino que está conectado a las herramientas del backend mediante el protocolo estándar \textbf{Model Context Protocol (MCP)} desarrollado en el directorio \texttt{mcp-server/}.
    \item \textbf{Orquestación segura de herramientas}: Cuando el usuario consulta al asistente cuestiones como «¿Cuánto dinero debo en el piso?» o «¿A quién le toca limpiar el baño esta semana?», el modelo de lenguaje selecciona e invoca dinámicamente las herramientas correspondientes (\texttt{getPendingDebts}, \texttt{getChoreSchedule}). Para preservar la privacidad entre hogares, la ejecución de cada herramienta exige un ticket de sesión temporal que acota estrictamente la consulta a los datos del piso al que pertenece el usuario.
\end{enumerate}

